\subsection{\texorpdfstring{Hermitian and \(\varepsilon\)-hermitian forms.}{Hermitian and epsilon-hermitian forms.}}

DEFINITION 1. --- Let \(\varepsilon\) be an element of the center of A. A sesquilinear form \(\Phi\) on E such that one has \(\Phi(x, y) = \varepsilon \overline{\Phi(y, x)}\) for all \(x\) and \(y\) in E is called a form \(\varepsilon\) -hermitian. A 1-hermitian form (resp. \((-1)\) -hermitian) is said to be hermitian (resp. antihermitian).

When J is the identity (which implies that A is commutative), a Hermitian (resp. antihermitian) form (for J) is none other than a symmetric (resp. antisymmetric) bilinear form (chap.~III, § 5, no. 1, def. 2). Recall that an alternating bilinear form (chap.~III, § 5, no. 2, def. 4) is antisymmetric; the converse is true if, in A, the relation \(2a = 0\) implies \(a = 0\) .

The orthogonality relation (§ 1, no. 3) with respect to an ε-hermitian form is evidently symmetric (cf.~exerc. 1).

If \(\alpha\) is an invertible element of A, the map T: \(\lambda \to \alpha^{-1}\overline{\lambda}\alpha\) is an antiautomorphism of A, and one readily verifies that the form \(\Phi\alpha\) is sesquilinear with respect to T. If, moreover, we have \(\alpha = \overline{\alpha}\) then T is involutive, and, if \(\Phi\) is \(\varepsilon\) -hermitian, so is \(\Phi\alpha\) ; indeed we have

\[
(\lambda^ {\mathrm{T}}) ^ {\mathrm{T}} = \alpha^ {- 1} (\overline {{\alpha^ {- 1} \lambda \alpha}}) \alpha = \alpha^ {- 1} \overline {{\alpha}} \lambda \overline {{\alpha}} ^ {- 1} \alpha = \lambda
\]

\[
\Phi (y, x) \alpha = \varepsilon \overline {{\Phi (x , y)}} \alpha = \varepsilon (\Phi (x, y) \alpha) ^ {\mathrm{T}}.
\]

In particular, when A is a field, the elements \(\alpha\) of the center of A such that \(\bar{\alpha} = \alpha\) form a subfield K of A, and the \(\varepsilon\) -hermitian forms on E (for J) form a vector space over K.

Remarks. --- 1) If \(\Phi\) is a form \(\varepsilon\) -hermitian on E, we have \(\Phi(x, y) = \varepsilon \Phi(x, y) \bar{\varepsilon}\) for all \(x, y\) in E. Therefore, if \(\Phi\) takes invertible values, we have \(\varepsilon \bar{\varepsilon} = 1\) .

\begin{enumerate}
\def\labelenumi{\arabic{enumi})}
\setcounter{enumi}{1}
\tightlist
\item
  If there exists an invertible element i of the center of A such that \(\bar{i}=i\varepsilon\) , then, for \(\Phi\) to be \(\varepsilon\) -hermitian, it is necessary and sufficient that \(i\Phi\) be Hermitian.
\end{enumerate}

The map \((y, x) \to \overline{\Phi(x, y)}\) being sesquilinear for J, for \(\Phi\) to be \(\varepsilon\) -hermitian, it is necessary and sufficient that one has \(\Phi(y, x) = \varepsilon \overline{\Phi(x, y)}\) when \(x\) and \(y\) range over a system of generators of E. In particular, if E admits a finite basis \((e_i)_{1 \leqslant i \leqslant n}\) for a sesquilinear form \(\Phi\) on E to be \(\varepsilon\) -hermitian, it is necessary and sufficient that its matrix \(R = (\rho_{ij}) = (\Phi(e_i, e_j))\) satisfies the relations \(\rho_{ji} = \varepsilon \overline{\rho_{ij}}\) for all \(i, j\) , that is to say \(^t R = \varepsilon \overline{R}\) ; a matrix \(R\) possessing this property is said to be \(\varepsilon\) -Hermitian. When \(\varepsilon = 1\) (resp. -1) one says that \(R\) is Hermitian (resp. antihermitian) with respect to the anti-automorphism J. When J is the identity (hence A is commutative), a Hermitian (resp. antihermitian) matrix \(R\) is such that \(^t R = R\) (resp. \(^t R = -R\) ); we then say that \(R\) is a symmetric (resp. antisymmetric) matrix. For \(\Phi\) to be an alternating form, it is necessary and sufficient that its matrix be antisymmetric and, in addition, that the diagonal terms of R all be zero; a matrix possessing these properties is said to be alternating.

Let \(\Phi\) a sesquilinear form on E, and let \(s_{\Phi}\) and \(d_{\Phi}\) be the maps from E into E* associated with \(\Phi\) to the left and to the right (§ 1, no. 6). For \(\Phi\) to be \(\varepsilon\) -hermitian, it is necessary and sufficient that \(\langle x, s_{\Phi}(y) \rangle = \bar{\varepsilon} \langle x, d_{\Phi}(y) \rangle\) for all elements \(x, y\) of E, hence that \(s_{\Phi} = \bar{\varepsilon} d_{\Phi}\) , or equivalently that \(\langle x, d_{\Phi}(y) \rangle = \varepsilon \langle x, s_{\Phi}(y) \rangle\) , hence that \(d_{\Phi} = \varepsilon s_{\Phi}\) .

Let \(\Phi\) be a form \(\varepsilon\) -hermitian such that the map \(d_{\Phi}\) from E into E* associated on the right to \(\Phi\) is bijective. For every endomorphism \(u\) of E we then have

\[
u ^ {* *} = \varepsilon \bar {\varepsilon} u.\tag{1}
\]

Indeed, whatever the elements \(x\) and \(y\) of \(\mathbf{E}\) , one has

\[
\begin{array}{r l} \Phi (x, u ^ {* *} (y)) & = \Phi (u ^ {*} (x), y) = \varepsilon \overline {{\Phi (y , u ^ {*} (x))}} = \varepsilon \overline {{\Phi (u (y) , x)}} \\ & = \varepsilon \Phi (x, u (y)) \bar {\varepsilon} = \Phi (x, \varepsilon \bar {\varepsilon} u (y)) \end{array}
\]

therefore \(u^{**}(x) = \varepsilon \bar{\varepsilon} u(x)\) since \(\Phi\) is nondegenerate.

If \(\Phi\) is a form \(\varepsilon\) -hermitian such that the maps \(s_{\Phi}\) and \(d_{\Phi}\) are bijective, then the inverse form \(\hat{\Phi}\) of \(\Phi\) ( \(\S 1\) , no. 7) is a \(\bar{\varepsilon}\) -hermitian form. Indeed, setting \(s = s_{\Phi}\) , \(d = d_{\Phi}\) for brevity, one deduces from \(d = \varepsilon s\) that \(s^{-1} = \bar{\varepsilon} d^{-1}\) , \(s\) being semilinear. Consequently, whatever \(u\) , \(\nu\) in E, one has

\[
\hat {\Phi} (u, \nu) = \Phi (s ^ {- 1} (u), d ^ {- 1} (\nu)) = \bar {\varepsilon} \Phi (d ^ {- 1} (u), d ^ {- 1} (\nu)),
\]

whence

\[
\widehat {\Phi} (\nu , u) = \overline {{\varepsilon \varepsilon}} \overline {{\Phi (d ^ {- 1} (u) , d ^ {- 1} (\nu))}} = \overline {{\varepsilon}} \widehat {\Phi} (u, \nu),
\]

since ε is in the center of A.

Finally, when the ring A is commutative, the canonical extensions of a \(\varepsilon\) -hermitian form \(\Phi\) to the tensor and exterior powers \(\bigotimes^{p}\mathbf{E}\) and \(\bigwedge^{p}\mathbf{E}\) of E are \(\varepsilon^{p}\) -hermitian forms, as follows immediately from formulas (35) and (37) of § 1, no. 9.

